\section{Methods}
\label{sec:methods}
\subsection{Surge definitions}
Let $y_t$ be the cases in week $t$ and $h$ the forecast horizon.
\begin{itemize}
  \item \textbf{Threshold (A):} $S^{A}_{t,h} = \mathbb{1}\{y_{t+h} > \mu_t + k\,\sigma_t\}$,
        where $\mu_t,\sigma_t$ are the mean and standard deviation of the previous $W=52$ weeks.
  \item \textbf{Growth (B):} $S^{B}_{t,h} = \mathbb{1}\{y_{t+h} \ge (1+p)\,y_t \;\wedge\; y_{t+h}\ge m\}$.
\end{itemize}

\subsection{Predictors}
All predictors are known at week $t$: lagged log cases, rolling means, recent growth,
seasonal terms $\sin/\cos$ of week-of-year, and weather lagged 0--12 weeks.

\subsection{Models}
Logistic regression (unpenalised and $\ell_1$), a depth-3 classification tree,
random forest, gradient boosting, and Poisson count regressions whose forecasts are
converted to alarms with the same rule that defines the target.

\subsection{Validation and model selection}
Rolling-origin evaluation: each year from 2017 is a test fold, trained on all earlier
weeks, with the last $h$ weeks purged to avoid target leakage. Metrics: ROC-AUC,
PR-AUC (preferred under class imbalance), Brier score, precision, recall and false alarm rate.
